\chapter{TỔNG KẾT VÀ HƯỚNG PHÁT TRIỂN} % Tên của chương

\label{Chapter5} % Để trích dẫn chương này ở chỗ nào đó trong bài, hãy sử dụng lệnh \ref{Chapter1} 

%----------------------------------------------------------------------------------------

% Định nghĩa một số lệnh cần thiết để điều chỉnh định dạng cho một số nội dung nhất định trong bài
\renewcommand{\keyword}[1]{\textbf{#1}}
\renewcommand{\tabhead}[1]{\textbf{#1}}
\renewcommand{\code}[1]{\texttt{#1}}
\renewcommand{\file}[1]{\texttt{\bfseries#1}}
\renewcommand{\option}[1]{\texttt{\itshape#1}}
\setcounter{secnumdepth}{3} 

\section{Tổng kết và đánh giá hệ thống luận chứng}

Trải qua các giai đoạn kiểm thử độc lập và thực nghiệm hợp nhất, nghiên cứu đã chuẩn hóa thành công quy trình vận hành đồng bộ của hệ thống xử lý luồng dữ liệu thời gian thực. Sự phối hợp chặt chẽ giữa tầng nhận thức đối tượng mục tiêu (YOLO26) \cite{Reference3}, tầng phân tích động lực học hành vi (ByteTrack) \cite{Reference12}, tầng suy diễn hình học không gian (MiDAS) \cite{Reference11} và tầng phản hồi ngôn ngữ tự nhiên (TTSEngine) đã thiết lập một kiến trúc xử lý khép kín hoàn chỉnh (\textit{End-to-End Pipeline}). Sự liên kết này không chỉ đảm bảo tính an toàn tuyệt đối mà còn tối ưu hóa khả năng tương tác liên tục, hỗ trợ người khiếm thị làm chủ không gian di chuyển.

Chu kỳ xử lý và lan truyền thông tin của hệ thống đối với mỗi khung hình dữ liệu đầu vào từ cảm biến quang học được mô hình hóa toán học theo chuỗi biến đổi trạng thái tương tác liên tục. Khung hình thu nhận trực tiếp từ cảm biến camera sẽ được lõi mạng học sâu YOLO26 trích xuất các hộp bao vị trí và nhãn phân lớp ngữ nghĩa vật cản như lối đi, đường gạch nổi, hay người đi bộ \cite{Reference3}. Ngay sau đó, trạng thái động lực học bao gồm mã định danh nhất quán, vectơ vận tốc pixel phẳng và xu hướng hướng dịch chuyển tương đối được đồng bộ hóa liên tục bởi bộ theo dõi ByteTrack thông qua liên kết lân cận Euclide \cite{Reference12}. Dữ liệu này kết hợp với điểm số chiều sâu đại diện chuẩn hóa của thực thể, được trích xuất bằng hàm thống kê phân vị thứ 80 trên ma trận nội suy đa khối của MiDAS hoặc cơ chế dự phòng tỷ lệ diện tích hình học, tạo thành chuỗi ngữ nghĩa văn bản mang đầy đủ thông tin ngữ cảnh biến đổi để chuyển giao trực tiếp vào bộ máy phát lệnh thoại \cite{Reference11}.

Điểm mấu chốt quyết định tính khả thi vượt trội của nghiên cứu khi triển khai trên các thiết bị phần cứng giới hạn năng lực tính toán nằm ở sự cân bằng chiến lược giữa huấn luyện đặc hiệu và tri thức chuyển giao (\textit{Transfer Learning}). Hệ thống không phân rã tài nguyên phần cứng cho các mạng nơ-ron đa nhiệm khổng lồ mà tập trung tối đa cho việc huấn luyện chuyên sâu mô hình gọn nhẹ YOLO26 nhằm đạt độ nhạy bén phân lớp tuyệt đối trên tập dữ liệu chướng ngại vật đặc thù của hạ tầng giao thông đô thị Việt Nam \cite{Reference3}. Đối với chiều không gian, hệ thống giải phóng bộ đệm lưu trữ bằng cách gọi suy diễn \textit{zero-shot} trực tiếp cấu trúc mạng \textit{MiDAS\_small} tiền huấn luyện, tận dụng khả năng bóc tách biên độ chiều sâu diện rộng mà không cần tiêu tốn chu kỳ lan truyền ngược (\textit{Backpropagation}) \cite{Reference11}.

Bên cạnh đó, cơ chế đồng bộ hóa luồng song song và triệt tiêu độ trễ tích lũy đóng vai trò cốt lõi trong việc duy trì hiệu năng thời gian thực. Sự kết hợp giữa bộ lọc khoảng cách Euclide tuyến tính ở tầng bám vết với giải pháp mã hóa ký số MD5 bộ nhớ đệm của TTSEngine đã triệt tiêu các phép toán ma trận phức tạp lặp lại. Tiến trình tổng hợp âm thanh được đẩy hoàn toàn vào các luồng phụ bất đồng bộ (\textit{Multithreading}) song song với luồng xử lý ảnh gốc. Đặc biệt, lõi dự phòng ngoại tuyến \textit{pyttsx3} phối hợp cùng hệ thống chuyển đổi tần số \textit{FFmpeg} đóng vai trò như một bộ đệm an toàn, giữ cho tốc độ xử lý khung hình luôn duy trì ở ngưỡng cao ổn định, loại bỏ hoàn toàn hiện tượng nghẽn luồng lệnh thoại khi thiết bị di chuyển qua các vùng mất kết nối mạng Internet \cite{Reference13}.

Tóm lại, tổng thể kiến trúc đề xuất là một phương pháp lai tối ưu, kết hợp hài hòa giữa mô hình học máy định hướng mục tiêu và các thuật toán hình học toán học, quản lý dữ liệu truyền thống. Sự liên kết này là minh chứng khoa học cho thấy giải pháp nghiên cứu đã vượt qua giới hạn của một mô hình lý thuyết thuần túy để trở thành một hệ thống nhúng hỗ trợ di chuyển có tính thực tiễn cao, vận hành bền bỉ, nhạy bén và là một công cụ trợ giúp đắc lực, đáng tin cậy cho cộng đồng người khiếm thị trong đời sống hàng ngày.

\section{Hạn chế và hướng phát triển}
 
Mặc dù hệ thống đề xuất đã đạt được những kết quả thực nghiệm khả quan trong việc nhận biết vật cản và hỗ trợ người khiếm thị di chuyển thời gian thực, nghiên cứu vẫn tồn tại một số hạn chế nhất định cần được nhìn nhận khách quan dưới góc độ kỹ thuật hệ thống. Trước hết, mô hình nhận diện chủ đạo YOLO26 tuy vận hành xuất sắc trên các nhóm vật cản tĩnh vỉa hè và hạ tầng giao thông cốt lõi, nhưng lại bộc lộ sự suy giảm hiệu năng đối với các phương tiện giao thông có kích thước lớn và cấu trúc hình học phức tạp như xe tải \cite{Reference3}. Hiện tượng che khuất cục bộ hoặc chồng lấn không gian giữa các thực thể chuyển động trong môi trường đô thị chật hẹp vẫn gây ra các sai số ngẫu nhiên cho bộ theo dõi ByteTrack, dẫn đến hiện tượng nhảy mã định danh hoặc trượt mục tiêu trong các khung hình chuyển tiếp ngắn \cite{Reference12}. Ngoài ra, việc phụ thuộc vào mô hình suy diễn ước tính chiều sâu đơn mục \textit{MiDAS\_small} đôi khi gặp khó khăn trong việc bóc tách chính xác ranh giới của các vật thể có bề mặt trong suốt như hệ thống cửa kính, hoặc các vật cản có hệ số phản xạ ánh sáng quá cao dưới điều kiện thời tiết cực đoan \cite{Reference11}. Cuối cùng, hệ thống phản hồi ngôn ngữ tự nhiên TTSEngine mặc dù đã tích hợp giải pháp đa luồng và bộ đệm cục bộ, nhưng cơ chế phát lệnh thoại tuần tự vẫn có thể gây ra hiện tượng quá tải thông tin âm thanh khi người dùng di chuyển vào các khu vực ngã tư đông đúc, nơi có quá nhiều chướng ngại vật xuất hiện đồng thời trong cùng một thời điểm tính toán \cite{Reference13}.

Từ những thách thức và hạn chế nêu trên, nghiên cứu định hướng các giải pháp phát triển và mở rộng hệ thống trong tương lai nhằm tối ưu hóa toàn diện sản phẩm thực tế. Về mặt nhận thức thị giác máy tính, hướng đi tiếp theo sẽ tập trung vào việc làm phong phú tập dữ liệu huấn luyện nội bộ bằng cách bổ sung đa dạng các trạng thái biến đổi góc nhìn của nhóm xe tải và phương tiện đặc thù, kết hợp ứng dụng các kỹ thuật tăng cường dữ liệu (\textit{Data Augmentation}) chuyên sâu để nâng cao chỉ số Recall tổng thể. Đối với tầng bám vết và ước tính không gian, nghiên cứu hướng tới việc tích hợp thêm các bộ lọc trạng thái Kalman thích nghi hoặc mô hình hình học dự phòng động, cho phép hệ thống tự động hiệu chuẩn khoảng cách dựa trên cả sự thay đổi của tiêu cự camera và độ dốc của mặt phẳng địa hình di chuyển. Để nâng cao độ tin cậy trong các kịch bản môi trường phức tạp, giải pháp tích hợp thêm cảm biến dòng phụ như cảm biến siêu âm hoặc bộ thu nhận tín hiệu định vị vệ tinh sẽ được cân nhắc triển khai nhằm tạo nên một cơ chế hợp nhất dữ liệu đa cảm biến (\textit{Sensor Fusion}) mạnh mẽ. Đặc biệt, đối với module phản hồi âm thanh, hệ thống sẽ được nâng cấp một thuật toán phân cấp ưu tiên động (\textit{Dynamic Priority Scheduling}) để lọc bỏ các thông báo không khẩn cấp, chỉ ưu tiên phát ngôn các cảnh báo có chỉ số nguy hiểm cao nhất theo thời gian thực, đồng thời tích hợp công nghệ âm thanh vòm 3D (\textit{Binaural Audio}) giúp người khiếm thị không chỉ nhận biết được khoảng cách mà còn định vị được chính xác phương hướng của vật cản trong không gian ba chiều.